\section{Introduction}\label{sec:intro}

The Birch--Swinnerton-Dyer formula relates the first nonzero Taylor
coefficient of an elliptic curve's $L$-function to its periods, rational
points, local components, and Tate--Shafarevich group. Knowing the order
of vanishing and the finiteness of the Tate--Shafarevich group does not
determine that coefficient. The problem addressed here is the remaining
integral comparison when a full Selmer group has corank zero or one.

\subsection{Statement and normalization}

Let $E/\Q$ be an elliptic curve. For a prime $p$, let
$\Sel_{p^\infty}(E/\Q)$ denote the full $p$-power Selmer group, defined
using the local Kummer images at every place, and write
\[
 s_p(E)=\corank_{\Z_p}\Sel_{p^\infty}(E/\Q),\qquad
 a(E)=\ord_{s=1}L(E,s),\qquad r(E)=\rank E(\Q).
\]
Here $L(E,s)$ is the uncompleted Hasse--Weil $L$-function with all its
finite Euler factors. Its center is $s=1$. We write
\[
 \Sha(E/\Q)=\ker\left(H^1(\Q,E)\longrightarrow
                          \prod_v H^1(\Q_v,E)\right).
\]

Choose a global minimal Weierstrass equation and its Neron differential
$\omega_E$. Our real period and finite component factors are
\[
 \Omega_E=\int_{E(\R)}|\omega_E|,
 \qquad c_\ell(E)=[E(\Q_\ell):E_0(\Q_\ell)].
\]
Thus the real period includes both connected components when there are
two. No separate real-component factor is used.

We fix the height convention as well. If $x(P)=a/b$ in lowest terms,
$b>0$, put $h_x(P)=\log\max\{|a|,b\}$ and $h_x(O)=0$. Define
\begin{equation}\label{eq:height-convention}
 H(P)=\lim_{n\to\infty}4^{-n}h_x([2^n]P),\qquad
 B(P,Q)=\frac{H(P+Q)-H(P)-H(Q)}2.
\end{equation}
For a basis $P_1,\ldots,P_r$ of the full free group
$E(\Q)/E(\Q)_{\rm tors}$, set
$\Reg_E=\det(B(P_i,P_j))$; in rank zero set $\Reg_E=1$.
In particular $B(P,P)=H(P)$. An index-$d$ sublattice has pairing
determinant $d^2\Reg_E$.

\begin{theorem}\label{thm:main}
Let $E/\Q$ be an elliptic curve and let $q$ be any prime. If
$s_q(E)\in\{0,1\}$, then
\[
 r(E)=a(E)=s_q(E),\qquad \#\Sha(E/\Q)<\infty,
\]
and, with $r=r(E)$ and the preceding normalizations,
\begin{equation}\label{eq:bsd-main}
 \frac{L^{(r)}(E,1)}{r!}
 =\frac{\Omega_E\Reg_E\,\#\Sha(E/\Q)
                  \prod_{\ell\text{ finite}}c_\ell(E)}
              {\#E(\Q)_{\rm tors}^{2}}.
\end{equation}
There are no additional hypotheses on reduction, rational torsion,
isogenies, complex multiplication, or residual Galois representations.
\end{theorem}

The rank and finiteness assertions are the unrestricted low-corank
converse of \cite{C}; the two-primary leading-term equality is supplied
by \cite{B}. We state these inputs precisely in
Section~\ref{sec:classical}. The new assertion proved here is the exact
valuation of the leading-term formula at every odd prime, independently
of the prime $q$ in the hypothesis. The low-Selmer-corank hypothesis is
essential to the stated scope: it is not replaced here by a hypothesis
on Mordell--Weil rank alone.

\subsection{Historical context and the integral problem}

Birch and Swinnerton-Dyer formulated their conjectures from the
relationship between the arithmetic of rational points and the behavior
of the Euler product \cite{BirchSwinnertonDyer}. Tate's account
\cite{TateBSD} placed the rank prediction and the leading-term formula
in the broader arithmetic of abelian varieties. The leading coefficient
contains information that the rank alone cannot see: finite local
components, torsion, the index of a point lattice, and the order of
\(\Sha\) must be compared simultaneously. Cassels's isogeny formula
\cite{Cassels} explains why this combined expression, rather than any
individual factor, is the natural arithmetic invariant.

The modularity theorems of Wiles, Taylor--Wiles, and
Breuil--Conrad--Diamond--Taylor
\cite{Wiles,TaylorWiles,BCDT} supply the analytic continuation and
functional equation of the \(L\)-function of every elliptic curve over
\(\Q\). Gross--Zagier's height formula \cite{GrossZagier} identifies
a derivative of a quadratic base-change \(L\)-function with the height
of a Heegner point. Kolyvagin's Euler-system argument
\cite{Kolyvagin} then gives rank equality and finiteness of the whole
Tate--Shafarevich group in analytic ranks zero and one. These results
are indispensable here, but their rank and finiteness conclusions do
not by themselves identify every prime valuation of the leading term.

Iwasawa theory provides a route to those integral refinements.
Rubin's main-conjecture theorem for imaginary quadratic fields
\cite{Rubin} gives leading-term consequences for CM elliptic curves.
Kato's zeta elements and explicit reciprocity law \cite{Kato}
connect modular \(L\)-values with Galois cohomology. In the non-CM
setting, the work of Skinner--Urban \cite{SkinnerUrban} yields
ordinary rank-zero leading-term results under residual and
ramification hypotheses. Jetchev--Skinner--Wan
\cite{JetchevSkinnerWan} prove rank-one formulas for semistable
curves at good primes under irreducibility hypotheses, including
ordinary and supersingular settings in their stated ranges.
These theorems make precise integral comparisons, with their own
local and residual assumptions; they are not used here as
unrestricted odd-primary formulas.

For CM curves, Burungale--Flach \cite{BurungaleFlach} prove the full
rank-zero BSD formula for their abelian torsion-field class.
Their rational descent applies to CM elliptic curves over \(\Q\).
That theorem supplies the rank-zero CM input in our final step.
The unrestricted low-corank converse \cite{C} and the two-primary
formula \cite{B} are distinct companion inputs. The auxiliary
twist-density theorem \cite{Goldfeld}, whose proof uses Smith's
Selmer-distribution results \cite{Smith}, permits simultaneous
choices of the two imaginary quadratic fields. We state the
scope of each input where it is used.

The cohomological language of the proof belongs to the arithmetic
duality and Selmer-complex framework
\cite{BlochKato,Milne,Nekovar}. In particular, exact local
orthogonality and determinant lines organize the local factors;
they do not remove the need to compute the integral lattices.
The moving-prime constructions use the Galois-image principles
of Serre and Bogomolov \cite{Serre,Bogomolov}, while analytic
primitivity uses the elliptic tame-stalk independence proved below
by the Chai--Hida rigidity and monodromy method
\cite{ChaiOrd,HidaMu}. The finite-model, switching, and theta
methods of \cite{B,C,Dv} provide the starting constructions.
The work below extends the relevant comparisons to residual
tests and to nonaugmentation theta tests with growing \(p\)-level,
and explains why these extensions give an exact center.

\subsection{The arithmetic discrepancy}

For an elliptic curve of analytic rank at most one, put
$r=r(E)=a(E)$ and define
\begin{equation}\label{eq:normalized-quotient}
 Q_E=\frac{L^{(r)}(E,1)\,\#E(\Q)_{\rm tors}^{2}}
            {r!\,\Omega_E\Reg_E\prod_\ell c_\ell(E)},
 \qquad
 X_p(E)=v_p(Q_E)-v_p(\#\Sha(E/\Q)),
\end{equation}
where $v_p(p)=1$. These quantities are defined: the analytic low-rank
theorem gives the rank and finiteness statements, and the two-primary
input gives $Q_E\in\Q_{>0}$. The latter applies because the Kummer exact
sequence gives $s_2(E)=r$ once the whole Tate--Shafarevich group is
finite.

Fix an odd prime $p$. The proof has two distinct outputs:
\begin{align}
 X_p(E)&\geq0 &&\text{if $E$ is non-CM and $a(E)\leq1$},
                                      \label{eq:intro-inequality}\\
 X_p(E)+X_p(E^D)&=0 &&\text{for a suitable imaginary quadratic twist},
                                      \label{eq:intro-pair}
\end{align}
where $a(E)+a(E^D)=1$. Applying the first inequality to both curves
in the second identity gives $X_p(E)=0$ in the non-CM case. The same
pair identity handles CM curves of analytic rank one after the CM
rank-zero formula is used. This separation is useful: the first
comparison is a one-sided integral statement, whereas the second
requires a unit, not just an integral quotient.

\subsection{Integral comparisons and the role of theta series}

The integral objects are finite free models of Galois cochain diagrams
over rings
\[
 R=\Z_p[[t,u_1,\ldots,u_m]].
\]
The variable $t$ is a genuine cyclotomic or anticyclotomic character
variable. The $u_i$ are limits of tame characters at varying auxiliary
primes. Every finite stage retains $T_pE$, the local conditions, and
the maps, pairings, and homotopies needed for derived specialization.
The tame variables make it possible to test and change the residual
cohomology while retaining an exact integral center.

The first comparison transfers Kato's modular-symbol construction into
these diagrams. Complementary local-condition comparisons prove that
its normalized determinant coordinate is integral. Exact calibration
against modular symbols in rank zero and a tame height formula in rank
one identifies its central valuation with $X_p(E)$. This proves
\eqref{eq:intro-inequality}; it does not require the normalized coordinate
to be a unit.

For the pair comparison, one works over
$K=\Q(\sqrt D)$, where $p$ splits as $w\bar w$. There are two determinant
series $L_w,L_{\bar w}$, obtained by imposing the strict (zero) local
condition at one of these places and the full (unrestricted) local
condition at the other. The analytic
series $B_w,B_{\bar w}$ are constructed from bounded CM measures.
After an integral coefficient descent, the relevant quotients are
\[
 U_w=B_w/L_w,\qquad U_{\bar w}=B_{\bar w}/L_{\bar w}.
\]
These formulas initially define elements of a fraction field. The
first task is to prove that they belong to $R$; the second is to prove
that they belong to $R^\times$.

Two different tests establish integrality and vertical primitivity.
Residual cohomology over $\Frac\F_p[[\mathbf u]]$ makes the determinants
regular in the tame variables modulo $(p,t)$. At sufficiently ramified
finite characters of $t$, local-condition comparisons give divisibility
after inverting $p$. Integral Weierstrass division then gives
$L_w\mid B_w$ over $R$. Independently, linear independence on the
ordinary modular curve proves $p\nmid B_w$. Thus $U_w$ is integral
and has no factor $p$. This does not yet exclude its other irreducible
factors.

Theta series are used for that remaining purpose. This part of the
proof allows a fixed loss of $p$-adic precision. At a characteristic-zero
divisor where an undepleted CM period vanishes, a theta-to-cusp
comparison produces a nonzero class satisfying a one-sided strict
condition, either for $E$ or for a fixed companion twist. The
comparison needs values at a nontrivial test character, not just the
augmentation. Its proof consequently allows increasing $p$-level and
proves a uniform bound for position-cut moments. A product-character
argument gives the required content bound at every interior test.

The theta implications make every possible quotient factor regular in
the tame variables. High-character association and the integral
character-division test then identify the two quotient ideals, including
their multiplicities. Adding
finitely many further tame variables then makes their possible
simultaneous strict jumps have codimension at least two. A quotient
divisor has codimension one, so no such divisor remains. Specializing
the added variables to zero is integral, and a unit remains a unit.
This last fact is what permits an exact central comparison; the
rationalized theta construction is not itself read as a formula for
the order of a finite $p$-primary group.

\begin{figure}[htbp]
\centering
\begin{tikzpicture}[
 box/.style={draw,align=center,text width=5.5cm,inner sep=6pt},
 >=Stealth,node distance=9mm]
 \node[box] (integral) {Integral character division\\
                         $U_w,U_{\bar w}\in R$, $p\nmid U_wU_{\bar w}$};
 \node[box,below=of integral] (theta) {Theta regularity and character division\\
                         $(U_w)=(U_{\bar w})$};
 \node[box,below=of theta] (remove) {New tame directions exclude\\
                         simultaneous jumps in codimension one};
 \node[box,below=of remove] (unit) {$U_w,U_{\bar w}\in R^\times$\\
                         Exact integral specialization at the center};
 \draw[->] (integral)--(theta);
 \draw[->] (theta)--(remove);
 \draw[->] (remove)--(unit);
\end{tikzpicture}
\caption{The pair comparison separates integral divisibility from
divisor removal. The theta implication is used on characteristic-zero
divisors; the final specialization takes place in the integral ring.}
\label{fig:pair-comparison}
\end{figure}

The reusable parts of this argument are the marked residual
obstruction calculation, the integral character-division test, and the
uniform position-cut and product-character estimates. Their roles are
different: the first controls determinant factors above $p$, the
second transfers divisibility from highly ramified characters, and the
third supplies the vanishing implication needed to remove horizontal
factors. We prove the changes needed for these applications rather
than inferring them from the conclusions of the companion papers.

\subsection{Organization}

Section~\ref{sec:classical} fixes the arithmetic inputs and the exact
height and period comparisons. Sections~\ref{sec:models}
and~\ref{sec:vertical} construct the integral diagrams and establish
the residual determinant tests. Section~\ref{sec:single} proves the
single-curve inequality, and Section~\ref{sec:pair} constructs the
integral pair quotients. Sections~\ref{sec:theta}
and~\ref{sec:extraction} prove the theta implication, including the
increasing-level and nonaugmentation estimates.
Section~\ref{sec:unit} removes the quotient divisors.
Section~\ref{sec:center} computes the central arithmetic volume and
completes the proof of Theorem~\ref{thm:main}.
